\documentclass[12pt]{article}
\usepackage{ficha}

\begin{document}

\ficha{Lugares Geométricos e Pontos Notáveis do Triângulo}

\instrucoes{Resolve os exercícios. Para as construções, usa régua, compasso e transferidor. As figuras são esquemáticas.}

\secao{Circunferência, Círculo e Mediatriz}{Descreve ou constrói o lugar geométrico pedido.}
\begin{exercicios}
  \item O lugar geométrico dos pontos do plano a 3\,cm de um ponto $O$.
  \item O lugar geométrico dos pontos do plano a uma distância de $O$ menor ou igual a 3\,cm.
  \item O lugar geométrico dos pontos do plano equidistantes de dois pontos $A$ e $B$.
  \item Desenha um segmento $[AB]$ com 6\,cm e constrói a sua mediatriz com régua e compasso.
  \item Um ponto $P$ pertence à mediatriz de $[AB]$ e $PA = 5$\,cm. Quanto mede $PB$? Justifica.
  \item A mediatriz de um segmento é perpendicular ao segmento? Passa pelo ponto médio?
  \item O lugar geométrico dos pontos do espaço a 2\,cm de um ponto $O$.
  \item O lugar geométrico dos pontos do plano a 2\,cm de uma reta $r$.
\end{exercicios}

\secao{Bissetriz de um Ângulo}{Responde a cada questão.}
\begin{exercicios}
  \item Define bissetriz de um ângulo como lugar geométrico.
  \item Desenha um ângulo de $70^\circ$ e constrói a sua bissetriz com régua e compasso.
  \item Um ponto $P$ pertence à bissetriz do ângulo $A\hat{O}B$ e está a 4\,cm do lado $\dot{O}A$. A que distância está do lado $\dot{O}B$?
  \item A bissetriz de um ângulo de $124^\circ$ divide-o em dois ângulos de que amplitude?
  \item Qual é o lugar geométrico dos pontos equidistantes de duas retas concorrentes?
\end{exercicios}

\secao{Pontos Notáveis do Triângulo}{Completa ou responde.}

\begin{center}
\begin{tikzpicture}[scale=0.9]
  \coordinate (A) at (0,0); \coordinate (B) at (6,0); \coordinate (C) at (2,4);
  \draw[thick] (A) node[below left] {$A$} -- (B) node[below right] {$B$} -- (C) node[above] {$C$} -- cycle;
\end{tikzpicture}
\end{center}

\begin{exercicios}
  \item O circuncentro é o ponto de interseção das \underline{\hspace{3cm}} do triângulo.
  \item O incentro é o ponto de interseção das \underline{\hspace{3cm}} do triângulo.
  \item O baricentro é o ponto de interseção das \underline{\hspace{3cm}} do triângulo.
  \item O ortocentro é o ponto de interseção das \underline{\hspace{3cm}} do triângulo.
  \item Qual dos pontos notáveis é o centro da circunferência circunscrita ao triângulo? Porquê?
  \item Qual dos pontos notáveis é o centro da circunferência inscrita no triângulo? Porquê?
  \item Copia o triângulo $ABC$ e constrói o seu circuncentro.
  \item Num triângulo retângulo, onde fica o circuncentro? E o ortocentro?
  \item Num triângulo obtusângulo, o circuncentro fica dentro ou fora do triângulo?
  \item Num triângulo equilátero, o que se pode dizer dos quatro pontos notáveis?
\end{exercicios}

\secao{Baricentro e Medianas}{Resolve.}
\begin{exercicios}
  \item O que é uma mediana de um triângulo?
  \item O baricentro divide cada mediana na razão $2:1$. Se uma mediana mede 12\,cm, a que distância do vértice fica o baricentro?
  \item Num referencial, os vértices de um triângulo são $(0,0)$, $(6,0)$ e $(0,9)$. Determina as coordenadas do baricentro (média das coordenadas dos vértices).
  \item Num referencial, os vértices de um triângulo são $(1,2)$, $(5,4)$ e $(3,-3)$. Determina o baricentro.
  \item Determina o ponto médio do segmento de extremos $(2, 6)$ e $(8, -2)$.
\end{exercicios}

\secao{Desafios}{Resolve os problemas.}
\begin{exercicios}
  \item Três aldeias $A$, $B$ e $C$ querem construir um poço à mesma distância das três. Como determinarias a localização do poço?
  \item Um jardim tem a forma de um triângulo. Quer-se instalar um aspersor circular que regue a maior zona possível sem molhar o exterior. Onde se deve colocar?
  \item Num triângulo retângulo de catetos 6\,cm e 8\,cm, qual é o raio da circunferência circunscrita?
  \item Constrói o lugar geométrico dos pontos que estão a menos de 3\,cm de $A$ e mais perto de $A$ do que de $B$, com $AB = 4$\,cm.
  \item Um tesouro está enterrado a igual distância de duas árvores e a 5\,m de um poço. Quantas posições possíveis podem existir? Justifica.
\end{exercicios}

\end{document}
